 \documentclass[conference]{IEEEtran}
\IEEEoverridecommandlockouts

\usepackage{cite}
\usepackage{amsmath,amssymb,amsfonts}
\usepackage{algorithmic}
\usepackage{algorithm}
\usepackage{graphicx}
\usepackage{textcomp}
\usepackage{xcolor}
\usepackage{booktabs}
\usepackage{multirow}
\usepackage{array}
\usepackage{amsthm}
\usepackage{pifont}
\usepackage{url}

\newtheorem{definition}{Definition}
\newtheorem{theorem}{Theorem}
\newtheorem{proposition}{Proposition}
\newtheorem{remark}{Remark}
\newtheorem{lemma}{Lemma}

\newcommand{\E}{\mathbb{E}}
\newcommand{\R}{\mathbb{R}}
% Indicator function rendered as bold 1 (universal, no extra package).
% On Overleaf you may instead use \usepackage{dsfont} + \mathds{1}
% or \usepackage{bbm} + \mathbbm{1} for a blackboard-bold look.
\newcommand{\indic}{\mathbf{1}}
\newcommand{\argmax}{\operatorname*{arg\,max}}
\newcommand{\JS}{\mathrm{JS}}
\newcommand{\Coh}{\mathrm{Coh}}

\def\BibTeX{{\rm B\kern-.05em{\sc i\kern-.025em b}\kern-.08em
    T\kern-.1667em\lower.7ex\hbox{E}\kern-.125emX}}

\begin{document}

\title{A Multi-Domain Cyber-Deception Bayesian Game\\
for Resilient IoT Systems}

\author{\IEEEauthorblockN{Alioum Abdoulaye}
\IEEEauthorblockA{\textit{Laboratoire DAVID} \\
\textit{Universit\'e Paris-Saclay}\\
Versailles, France \\
aliyoumabdoulaye@gmail.com}
\and
\IEEEauthorblockN{Yacine Benallouche}
\IEEEauthorblockA{\textit{Orange Research}\\ \textit{Orange} \\
Paris, France\\
yacine.benallouche@orange.com}
\and
\IEEEauthorblockN{Abdelhak Mourad Gueroui}
\IEEEauthorblockA{\textit{Laboratoire DAVID} \\
\textit{Universit\'e Paris-Saclay}\\
Versailles, France\\
mourad.gueroui@uvsq.fr}
\and
\IEEEauthorblockN{Ado Adamou Abba Ari}
\IEEEauthorblockA{\textit{Laboratoire DAVID} \\
\textit{Universit\'e Paris-Saclay}\\
Versailles, France\\
ado-adamou.abba-ari@uvsq.fr}
}

\maketitle

\pagestyle{plain}
\thispagestyle{plain}

\begin{abstract}
Internet of Things (IoT) infrastructures are increasingly targeted by
Advanced Persistent Threats (APTs) that progress methodically across the
perception, network, and application layers. Cyber-deception, which
strategically deploys decoys (honeypots) to mislead attackers, has emerged
as a proactive defense paradigm. However, existing game-theoretic deception
models suffer from three structural limitations: they treat each IoT layer in
isolation, they assume zero-sum interaction, and they ignore the resilience
requirements that govern critical IoT services. We introduce the
Multi-Domain Cyber-Deception Bayesian Game (MD-CDBG), a stochastic
Bayesian game with incomplete information that couples the three IoT layers
under formal resilience and semantic-coherence constraints. We define a novel
solution concept, the Resilience-constrained Bayesian Nash Equilibrium
(R-BNE), which embeds a hard resilience constraint at every time step, and we
solve it with a Bayesian Double Oracle algorithm whose inner oracle uses a
Linear Complementarity Program (LCP) to correctly handle the non-zero-sum
structure of the game. Through a Lagrangian decomposition, the three coupled
sub-games are solved in parallel while enforcing the budget and coherence
constraints, and a fallback isolation policy is triggered whenever resilience
approaches its threshold. Extensive simulations against three
baselines, including the recent GameSec~2025 double-oracle approach, across
three attacker profiles
(opportunistic, persistent, and state-sponsored) show that the MD-CDBG reduces
resilience-constraint violations by 65--72\% while inducing rational attacker
behavior, demonstrating that modeling deception, belief, and resilience jointly
is essential for protecting multi-layer IoT systems.
\end{abstract}

\begin{IEEEkeywords}
IoT security, cyber-deception, Bayesian game, Nash equilibrium, resilience,
double oracle, linear complementarity, honeypot, advanced persistent threat.
\end{IEEEkeywords}

% ══════════════════════════════════════════════════════════════
\section{Introduction}
% ══════════════════════════════════════════════════════════════

The proliferation of Internet of Things (IoT) devices in industrial,
medical, and urban infrastructures has dramatically expanded the attack
surface available to sophisticated adversaries. Billions of resource-constrained
sensors, gateways, and actuators now mediate critical services, and their
heterogeneity and scale make conventional perimeter defenses
insufficient~\cite{pawlick2019}. Among the threats they face, Advanced
Persistent Threats (APTs) are particularly dangerous: they operate methodically
over long periods, escalating privileges layer by layer through the canonical
three-tier IoT architecture, perception (sensors), network (gateways and
routers), and application (cloud services and decision engines)~\cite{wan2022}.

Traditional intrusion-detection systems are inherently reactive: they detect
attacks only after they manifest, by which point a persistent adversary may
have already compromised critical assets. Cyber-deception offers a
proactive alternative. By deploying decoys (honeypots) that mimic real assets,
a defender can mislead the attacker, waste its resources, and reveal its
intentions before genuine damage occurs~\cite{anwar2020,sayed2023}. The core
strength of cyber-deception is its proactive nature, transforming static
defenses into dynamic, unpredictable, and more resilient security
stances~\cite{asghar2025}. To be effective, however, deception must be
optimized through rigorous decision-making under uncertainty, for which game
theory provides a principled foundation~\cite{huang2019}.

Despite a rich literature, existing game-theoretic deception models exhibit
three structural gaps when applied to IoT. First, they typically
model a single domain or layer, ignoring the cross-layer coupling that
characterizes APT progression in IoT; a honeypot strategy optimal for the
network layer may be incoherent with the perception or application layer,
making decoys easy to unmask. Second, many adopt a zero-sum
formulation in which the attacker's gain is exactly the defender's loss; this
assumption fails for IoT, where energy consumption, deployment cost, and
resilience introduce defender objectives that are not mirrored by the attacker.
Third, and most critically, prior work largely ignores
resilience: the capacity of the system to maintain service continuity
under attack, which is the paramount concern for critical IoT
deployments~\cite{he2026}.

Contributions. This paper addresses these gaps with the following
contributions:
\begin{itemize}
\item We formulate the Multi-Domain Cyber-Deception Bayesian Game
(MD-CDBG), a stochastic Bayesian game coupling the three IoT layers under
formal resilience, semantic-coherence, and budget constraints
(Section~\ref{sec:model}).
\item We define the Resilience-constrained Bayesian Nash Equilibrium
(R-BNE), a solution concept that embeds a hard resilience constraint
$\Omega(t)\geq\Omega_{\min}$ at every time step (Section~\ref{sec:rbne}).
\item We design R-BNE-DO, a Bayesian Double Oracle algorithm whose
inner oracle solves the non-zero-sum game via a Linear Complementarity Program,
coordinated by a Lagrangian decomposition over the three layers, and we prove
its convergence (Section~\ref{sec:algo}).
\item We evaluate the MD-CDBG against three baselines, including the recent
GameSec~2025 double-oracle deception model~\cite{asghar2025}, across three
attacker profiles, demonstrating a 65--72\% reduction in resilience-constraint
violations across all scenarios and substantially more rational attacker
behavior, at a defender utility that matches or exceeds the strongest baseline
(Section~\ref{sec:results}).
\end{itemize}

The remainder of the paper is organized as follows.
Section~\ref{sec:related} reviews related work and positions our contribution.
Section~\ref{sec:model} formalizes the MD-CDBG. Section~\ref{sec:utilities}
defines the utility functions. Section~\ref{sec:rbne} introduces the resilience
function and the R-BNE. Section~\ref{sec:algo} presents the R-BNE-DO algorithm
and its convergence analysis. Section~\ref{sec:results} reports the
experimental evaluation. Section~\ref{sec:conclusion} concludes.

% ══════════════════════════════════════════════════════════════
\section{Related Work}
\label{sec:related}
% ══════════════════════════════════════════════════════════════

\subsection{Game-Theoretic Cyber-Deception}

Signaling games and Stackelberg games have been widely used to model the
strategic interaction between a defender deploying decoys and an attacker
probing a network~\cite{tambe2011}. Pawlick \textit{et al.}~\cite{pawlick2019}
provide a comprehensive taxonomy of defensive deception, categorizing
approaches by information structure and game type, and book-length treatments
establish the scientific foundations of cyber deception and of game theory and
machine learning for security~\cite{jajodia2016,kamhoua2021}. Huang and
Zhu~\cite{huang2019} model adaptive honeypot engagement as a semi-Markov
decision process solved by reinforcement learning. Two-layer signaling games
have been proposed to defend cyber-physical systems against attacks that
traverse both a cyber and a physical plane~\cite{kamdem2024}. These models
capture information asymmetry but generally treat a single domain and assume a
static type for the adversary, leaving cross-layer coupling and resilience
unaddressed.

\subsection{Double Oracle Methods for Large Games}

The Double Oracle (DO) meta-algorithm~\cite{mcmahan2003} computes equilibria in
large or infinite action spaces by iteratively expanding restricted strategy
sets, invoking a best-response oracle for each player until no profitable
deviation exists. DO and its variants have proven effective for
honeypot-allocation and attack-graph
games~\cite{nguemkam2024,ngo2021,durkota2015}. Most relevant to our work, Asghar
\textit{et al.}~\cite{asghar2025} model a dual-domain attack on cyber-physical
sensors as a non-zero-sum game solved by double oracle with a greedy,
submodular defender oracle, extending their earlier zero-sum multi-domain
deception formulation~\cite{asghar2024cns}. Crucially, although their game is
not zero-sum, they prove it is strategically equivalent to a zero-sum
game (in the sense of Moulin and Vial), which lets them solve each restricted
game with a single linear program and avoid the linear complementarity problem
(LCP) altogether. This equivalence, however, is a property of their particular
payoff structure: it does not survive once the defender's objective carries
resilience, energy-budget, and cross-layer coherence penalties that the attacker
does not share, as in resilient IoT. Their model also provides no resilience
guarantee. We show in Section~\ref{sec:model} that under these penalties the
game is not reducible to zero-sum, so the LCP becomes necessary rather
than optional.

\subsection{Bayesian Games and Belief Modeling}

Games with incomplete information, introduced by Harsanyi~\cite{harsanyi1967},
model uncertainty about an opponent's type through a prior belief updated by
Bayes' rule. In cyber-security, the attacker's type (e.g., its sophistication
or objective) is rarely known, motivating Bayesian formulations. Wan
\textit{et al.}~\cite{wan2022} extend this idea with hypergame theory, allowing
players to hold subjective and possibly inconsistent views of the same game,
and apply it to defensive deception against APTs. However, these approaches
maintain a belief over the attacker's type without coupling that belief to hard
operational constraints such as resilience or a global deployment budget.

\subsection{Deception in IoT and Wireless Sensor Networks}

A growing body of work targets deception in resource-constrained IoT and
wireless sensor networks (WSNs). Anwar \textit{et al.}~\cite{anwar2020} develop
a game-theoretic framework for dynamic cyber-deception in the Internet of
Battlefield Things. Sihomnou \textit{et al.}~\cite{sihomnou2025,sihomnou2025b}
study Bayesian and signaling games to defend IoT devices against energy-depletion
and battery-drain denial-of-service attacks, emphasizing the lightweight
constraints intrinsic to IoT. He \textit{et al.}~\cite{he2026} propose a deep
reinforcement learning approach to time-delay differential game deception for
real-time resource deployment. These works recognize IoT-specific constraints
but consider a single layer and do not formalize resilience as an equilibrium
constraint.

\subsection{Learning-Based and Detection Approaches}

Complementary to game theory, machine-learning detectors have been applied to
multi-domain intrusion detection, including deep autoencoder
models~\cite{musa2023} and graph neural networks that deceive network
reconnaissance~\cite{njanko2023}. Sayed \textit{et al.}~\cite{sayed2023}
combine game theory with learning to defend against zero-day attacks. Such
detectors learn attacker behavior reactively but lack the forward-looking
equilibrium reasoning and resilience guarantees that a game-theoretic
formulation provides; integrating such a learned detector as an additional
oracle within our framework is a promising direction
(Section~\ref{sec:conclusion}).

\subsection{Positioning}

To our knowledge, the MD-CDBG is the first model to (i) couple all three IoT
layers in a single Bayesian deception game, (ii) preserve the non-zero-sum
structure through an LCP-based oracle rather than a zero-sum proxy, and (iii)
embed resilience as a hard constraint within the equilibrium concept itself.
Table~\ref{tab:positioning} summarizes how the MD-CDBG differs from
representative prior work.

\begin{table}[t]
\centering
\caption{Positioning of the MD-CDBG against representative prior work.}
\label{tab:positioning}
\small
\begin{tabular}{lcccc}
\toprule
\textbf{Work} & \textbf{Multi-} & \textbf{Non-zero} & \textbf{Bayesian} &
\textbf{Resil.} \\
 & \textbf{layer} & \textbf{-sum} & \textbf{belief} & \textbf{constr.} \\
\midrule
Pawlick~\cite{pawlick2019}   & \ding{55} & \checkmark & \checkmark & \ding{55} \\
Huang~\cite{huang2019}       & \ding{55} & \checkmark & \ding{55} & \ding{55} \\
Wan~\cite{wan2022}           & \ding{55} & \checkmark & \checkmark & \ding{55} \\
Anwar~\cite{anwar2020}       & \ding{55} & \checkmark & \ding{55} & \ding{55} \\
Asghar~\cite{asghar2025}     & \checkmark & \ding{55} & \ding{55} & \ding{55} \\
Nguemkam~\cite{nguemkam2024} & \ding{55} & \checkmark & \checkmark & \ding{55} \\
\textbf{MD-CDBG (ours)}      & \checkmark & \checkmark & \checkmark & \checkmark \\
\bottomrule
\end{tabular}

\vspace{2pt}
{\footnotesize The ``non-zero-sum'' column indicates whether the equilibrium is
computed on the genuine non-zero-sum game. Asghar~\cite{asghar2025} formulate a
non-zero-sum game but reduce it to a strategically equivalent zero-sum game;
MD-CDBG instead solves the non-zero-sum game directly via the LCP.\par}
\end{table}

% ══════════════════════════════════════════════════════════════
\section{System Model: MD-CDBG}
\label{sec:model}
% ══════════════════════════════════════════════════════════════

\subsection{Game Definition}

We model the defender--attacker interaction across the three IoT layers as a
multi-domain stochastic Bayesian game:
\begin{equation}
\mathcal{G} = \langle \mathcal{G}_1, \mathcal{G}_2, \mathcal{G}_3,
\mathcal{C}, \Omega \rangle,
\end{equation}
where $\mathcal{G}_k$ is the sub-game at layer $k\in\{1,2,3\}$ (perception,
network, application), $\mathcal{C}=\{C_1,C_2,C_3\}$ is the set of coupling
constraints linking the sub-games, and $\Omega$ is the system resilience
function.

Each sub-game is a tuple
\begin{equation}
\mathcal{G}_k = \langle \mathcal{N}_k, \hat{\mathcal{N}}_k,
\mathcal{A}_k^D, \mathcal{A}_k^A, \Theta_A, u_k^D, u_k^A \rangle,
\end{equation}
comprising $|\mathcal{N}_k|$ real nodes and $|\hat{\mathcal{N}}_k|$ decoy nodes,
defender and attacker action sets $\mathcal{A}_k^D, \mathcal{A}_k^A$, a set of
attacker types $\Theta_A = \{\theta_{\mathrm{opp}}, \theta_{\mathrm{pers}},
\theta_{\mathrm{state}}\}$ corresponding to opportunistic, persistent, and
state-sponsored adversaries, and layer-level utilities $u_k^D, u_k^A$. The
defender chooses a mixed strategy $\sigma_k^D\in\Delta(\mathcal{A}_k^D)$; the
attacker of type $\theta_A$ chooses $\sigma_k^A(\cdot|\theta_A)
\in\Delta(\mathcal{A}_k^A)$. Figure~\ref{fig:arch} illustrates the overall
architecture, and Table~\ref{tab:notation} summarizes the main notation.

\begin{figure}[t]
\centering
\includegraphics[width=\columnwidth]{fig0.png}
\caption{Architecture of the MD-CDBG. The attacker progresses upward through
the three layers ($C_1$); decoys at adjacent layers must remain coherent
($C_2$); the defender maintains a per-layer Bayesian belief $\mu_k^t(\theta_A)$
over the attacker type, under a global budget ($C_3$) and a resilience
constraint $\Omega(t)\geq\Omega_{\min}$.}
\label{fig:arch}
\end{figure}

\begin{table}[t]
\centering
\caption{Main notation.}
\label{tab:notation}
\small
\begin{tabular}{ll}
\toprule
\textbf{Symbol} & \textbf{Meaning} \\
\midrule
$\mathcal{G}_k$ & Sub-game at layer $k$ (perception/network/app.) \\
$\mathcal{N}_k,\hat{\mathcal{N}}_k$ & Real and decoy nodes at layer $k$ \\
$\Theta_A$ & Attacker types (opp./pers./state) \\
$\mathcal{A}_k^D,\mathcal{A}_k^A$ & Defender / attacker action sets at layer $k$ \\
$\Delta(\mathcal{A}_k^D),\Delta(\mathcal{A}_k^A)$ & Probability simplices over actions \\
$a_{k,i}^D,a_{k,i}^A$ & $i$-th defender / attacker action at layer $k$ \\
$a^A(t),a^D(t)$ & Joint attacker / defender action profiles \\
$a_k^A(t)$ & Observed attacker action at layer $k$, time $t$ \\
$\mathbf{s},\mathbf{s}_k(t)$ & Global / per-layer system state profile \\
$\mathcal{T}$ & State-transition function \\
$\theta_A^*$ & True (realized) attacker type \\
$\sigma_k^D,\sigma_k^A$ & Defender / attacker mixed strategies \\
$\mu_k^t(\theta_A)$ & Bayesian belief on type at layer $k$, time $t$ \\ 
$L_t$ & Likelihood-ratio process\\
$\delta_k(t)$ & Compromise ratio of layer $k$ \\
$f_1^k$ & Deception effectiveness \\
$f_3^k,f_4^k$ & Energy overhead and deployment cost \\
$q_k$ & Composite decoy quality \\
$q_k^{\mathrm{func}},q_k^{\mathrm{stat}},q_k^{\mathrm{temp}}$ & Functional / statistical / temporal realism \\
$r_{k,i}^{\mathrm{func}}$ & Intrinsic realism of action $a_{k,i}^D$ \\
$\omega_i$ & Decoy-quality component weights \\
$d_k,d_k^{\mathrm{eff}},d_k^{\mathrm{vuln}}$ & Detection capability (active/passive) \\
$\psi_k(\theta_A)$ & Type-specific detection efficiency \\
$\kappa$ & Passive-detection coefficient \\
$\Coh_k$ & Semantic coherence between layers $k,k{+}1$ \\
$\Omega_k,\Omega$ & Local / global resilience \\
$\mathrm{res}_k,\mathrm{abs}_k,\mathrm{rec}_k$ & Resistance / absorption / recovery terms \\
$w_1,w_2,w_3$ & Adaptive resilience-component weights \\
$\rho_k$ & Resistance threshold (layer $k$) \\
$\tau_k,T_k^{\max}$ & Per-node and max recovery time \\
$\Omega_{\min}$ & Resilience threshold (hard constraint) \\
$\Pi^D_\Omega$ & Admissible (feasible) defender policy set \\
$\sigma^{D*},\sigma^{A*}$ & R-BNE defender / attacker policies \\
$g_k^A,c_k^A,\rho_k^A$ & Attacker gain / cost / detection risk \\
$v_{k,i}^A,p_{k,i}^A$ & Asset value, success probability \\
$r_{k,i}^A,\lambda_{k,i}$ & Action resource cost, detection sensitivity \\
$e_{k,i}^D,c_{k,i}^D$ & Defender unit energy / deployment cost \\
$u_k^D,u_k^A$ & Local defender / attacker utility \\
$U^D,U^A$ & Global defender / attacker utility \\
$\mathcal{L}_k$ & Per-layer Lagrangian \\
$M_k^D,M_k^A$ & Defender / attacker payoff matrices (Bayes-DO) \\
$R_D,R_A$ & Restricted action sets (double oracle) \\
$\alpha_i,\gamma_i$ & Defender / attacker utility weights \\
$\beta_k,\beta_k^A,\beta_k^\Omega$ & Layer aggregation weights \\
$\gamma$ & Discount factor; \;$T$: time horizon \\
$\eta_t,\eta_0$ & Gradient step size (Robbins--Monro) \\
$\epsilon,\varepsilon$ & Stability constant; convergence tolerance \\
$\lambda$ & Lagrange multiplier (budget $C_3$) \\
$\nu_k$ & Lagrange multiplier (coherence $C_2$) \\
$\delta_{\min}$, $B$ & Coherence threshold / budget \\
\bottomrule
\end{tabular}
\end{table}

\subsection{Coupling Constraints}

The three sub-games are linked by three constraints that make the MD-CDBG a
genuinely multi-domain game rather than three independent games:
\begin{itemize}
\item ($C_1$) APT progression: layer $k$ is accessible only if at
least one node in layer $k-1$ is compromised. This enforces the sequential
escalation characteristic of APTs: an attacker must breach the perception layer
before the network layer, and the network before the application.
\item ($C_2$) Semantic coherence: decoys at adjacent layers must be
mutually consistent, $\Coh(\hat{\mathcal{N}}_k,\hat{\mathcal{N}}_{k+1}) \geq
\delta_{\min}$. An attacker that observes a network decoy inconsistent with the
perception decoys it has already seen can unmask the deception.
\item ($C_3$) Global budget: the total deployment cost across all
layers must not exceed a budget $B$, i.e. $\sum_k f_4^k \leq B$, reflecting the
finite operational resources of an IoT operator.
\end{itemize}

\subsection{State and Transition}

The system state $\mathbf{s}(t) = (\mathbf{s}_1(t), \mathbf{s}_2(t),
\mathbf{s}_3(t))$ records, for each layer, the set of compromised nodes. The
compromise ratio of layer $k$ is $\delta_k(t) = |\text{compromised}_k| /
|\mathcal{N}_k|$. At each step, the realized defender and attacker actions
induce a stochastic transition $\mathbf{s}(t{+}1) =
\mathcal{T}(\mathbf{s}(t), a^A(t), a^D(t))$: an attacker action targeting a real
node succeeds with a technical success probability $p_{k,i}^A$ unless it is
trapped by a decoy (probability $f_1^k$), while a defender isolation action can
recover a compromised node.

\subsection{Bayesian Belief}

Since the attacker's true type $\theta_A^*$ is unobservable, the defender
maintains a belief $\mu_k^t(\theta_A) \in \Delta(\Theta_A)$ per layer
$k$ and time $t$, initialized at the prior $\mu_k^0(\theta_A) = p(\theta_A)$ and
updated after each observed action $a_k^A(t)$ via Bayes' rule:
\begin{equation}
\mu_k^{t+1}(\theta_A) =
\frac{\sigma_k^A(a_k^A(t)\mid\theta_A)\,\mu_k^t(\theta_A)}
{\sum_{\theta'\in\Theta_A}\sigma_k^A(a_k^A(t)\mid\theta')\,\mu_k^t(\theta')}.
\label{eq:bayes}
\end{equation}

\begin{lemma}[Belief convergence]
\label{lem:belief}
Suppose the type-conditional strategies
$\{\sigma_k^A(\cdot\mid\theta_A)\}_{\theta_A\in\Theta_A}$ are pairwise distinct
(\emph{identifiability}) and layer $k$ is visited infinitely often. Then, under
the true type $\theta_A^*$, the belief $\mu_k^t(\theta_A^*)\to 1$ almost surely
as $t\to\infty$; that is, the belief concentrates on the true type. The result
follows because, for every $\theta\neq\theta_A^*$, the likelihood-ratio process
$L_t(\theta)=\mu_k^t(\theta)/\mu_k^t(\theta_A^*)$ is a non-negative martingale
under $\theta_A^*$ and hence converges almost surely; the strict positivity of
the Kullback--Leibler divergence between distinct type-conditional
distributions forces its limit to $0$, so all false types vanish and
$\mu_k^t(\theta_A^*)\to 1$.
\end{lemma}

% ══════════════════════════════════════════════════════════════
\section{Utility Functions}
\label{sec:utilities}
% ══════════════════════════════════════════════════════════════

\subsection{Decoy Quality and Detection}

The realism of a layer-$k$ decoy deployment is captured by a composite quality
$q_k(\sigma_k^D)\in[0,1]$ combining functional, statistical, and temporal
realism:
\begin{equation}
q_k(\sigma_k^D) = \omega_1 q_k^{\mathrm{func}} + \omega_2 q_k^{\mathrm{stat}}
+ \omega_3 q_k^{\mathrm{temp}},
\qquad \textstyle\sum_i \omega_i = 1,
\label{eq:quality}
\end{equation}
where the functional realism $q_k^{\mathrm{func}}(\sigma_k^D)=\sum_i
 \sigma_k^D(a_{k,i}^D)\,r_{k,i}^{\mathrm{func}}$ is the strategy-weighted
realism of the chosen decoy actions ($r_{k,i}^{\mathrm{func}}\in[0,1]$ being the
intrinsic realism of action $a_{k,i}^D$), and the statistical and temporal
components measure how closely decoy traffic matches that of real nodes.

The attacker's detection capability combines an active and a passive component
through a complementary-probability law:
\begin{equation}
d_k(\sigma_k^A,\sigma_k^D) =
1 - \bigl(1-d_k^{\mathrm{eff}}\bigr)\bigl(1-d_k^{\mathrm{vuln}}\bigr),
\label{eq:detection}
\end{equation}
where the active component scales the attacker's probing investment by
its type-specific efficiency,
\begin{equation}
d_k^{\mathrm{eff}}(\sigma_k^A) = \sigma_k^A(a_k^{\mathrm{probe}})
\cdot \psi_k(\theta_A),
\end{equation}
and the passive component reflects the intrinsic detectability of
imperfect decoys, $d_k^{\mathrm{vuln}}(\sigma_k^D) = \kappa\,(1-q_k)$ with
$\kappa\in[0,1]$. Equation~\eqref{eq:detection} expresses that a decoy escapes
detection only if it survives both the active probe and the passive
inspection, hence the product of the two survival probabilities.

\subsection{Deception Effectiveness}

The deception effectiveness at layer $k$ multiplies the decoy ratio, the decoy
quality, and the complement of the detection capability:
\begin{equation}
f_1^k(\sigma_k^D,\sigma_k^A) =
\frac{|\hat{\mathcal{N}}_k|}{|\mathcal{N}_k|+|\hat{\mathcal{N}}_k|}
\cdot q_k(\sigma_k^D)\cdot\bigl(1-d_k(\sigma_k^A,\sigma_k^D)\bigr).
\label{eq:f1}
\end{equation}
The three factors admit a clear reading: more decoys raise the probability that
the attacker targets one; higher quality reduces the chance the attacker
recognizes it; and lower detection capability increases the chance the decoy
goes unnoticed.

\subsection{Semantic Coherence}

The coherence constraint $C_2$ penalizes abrupt quality discontinuities between
adjacent layers, which an attacker can exploit to unmask the deception. We
measure coherence by the normalized similarity of the decoy quality at two
adjacent layers:
\begin{equation}
\Coh_k(\sigma_k^D,\sigma_{k+1}^D) =
1 - \frac{\bigl|\,q_k(\sigma_k^D) - q_{k+1}(\sigma_{k+1}^D)\,\bigr|}
{\max\!\bigl(q_k,\,q_{k+1},\,\epsilon\bigr)},
\label{eq:coherence}
\end{equation}
where $\epsilon>0$ avoids division by zero. $\Coh_k\in[0,1]$ equals $1$ when the
two layers deploy decoys of identical quality and decreases as their realism
diverges; constraint $C_2$ requires $\Coh_k\geq\delta_{\min}$ for all adjacent
pairs.

\subsection{Defender Utility}

\begin{definition}[Defender utility]\label{def:util_D}
The defender's utility at layer $k$ is a weighted linear combination of four
partially conflicting objectives:
\begin{align}
u_k^D&\!\left(\sigma_k^D,\sigma_k^A(\cdot|\theta_A)\right) =
\underbrace{\alpha_1 f_1^k}_{\text{deception}}
+ \underbrace{\alpha_2 \Omega_k(\sigma_k^D)}_{\text{resilience}} \notag\\
&\quad - \underbrace{\alpha_3 f_3^k(\sigma_k^D)}_{\text{energy}}
- \underbrace{\alpha_4 f_4^k(\sigma_k^D)}_{\text{cost}},
\label{eq:util_D_k}
\end{align}
with $\sum_i\alpha_i=1$, $\alpha_i\geq 0$. The energy overhead and operational
deployment cost are the strategy-weighted sums of the per-action unit costs:
\begin{equation}
f_3^k(\sigma_k^D) = \sum_i \sigma_k^D(a_{k,i}^D)\,e_{k,i}^D,
\qquad
f_4^k(\sigma_k^D) = \sum_i \sigma_k^D(a_{k,i}^D)\,c_{k,i}^D,
\label{eq:costs}
\end{equation}
where $e_{k,i}^D$ and $c_{k,i}^D$ are the unit energy and deployment costs of
defender action $a_{k,i}^D$ (both normalized to $[0,1]$).
\end{definition}

Because the true type is unobservable, the defender evaluates $u_k^D$ in
expectation over its current belief, giving the belief-weighted local utility
\begin{equation}
\bar{u}_k^D\!\left(\sigma_k^D(\mu_k^t),\mu_k^t\right) =
\E_{\theta_A\sim\mu_k^t}\!\left[u_k^D\!\left(\sigma_k^D,
\sigma_k^A(\cdot|\theta_A)\right)\right],
\end{equation}
and the global defender utility aggregates over layers:
\begin{equation}
U^D(\sigma^D(\mu^t),\mu^t) = \sum_{k=1}^3 \beta_k\,
\bar{u}_k^D\!\left(\sigma_k^D(\mu_k^t),\mu_k^t\right),
\label{eq:util_D_global}
\end{equation}
with $\sum_k\beta_k=1$, $\beta_k>0$ reflecting the relative criticality of each
layer.

\subsection{Attacker Utility}

The utility of attacker type $\theta_A$ at layer $k$ balances expected asset
gain $g_k^A$, resource cost $c_k^A$, and detection risk $\rho_k^A$:
\begin{equation}
u_k^A(\sigma_k^A,\sigma_k^D\mid\theta_A) =
\gamma_1 g_k^A - \gamma_2 c_k^A - \gamma_3 \rho_k^A,
\label{eq:util_A}
\end{equation}
with $\sum_i\gamma_i=1$. The three components are
\begin{align}
g_k^A &= \sum_i \sigma_k^A(a_{k,i}^A)\,v_{k,i}^A\,(1-f_1^k)\,p_{k,i}^A,
\label{eq:gain}\\
c_k^A &= \sum_i \sigma_k^A(a_{k,i}^A)\,r_{k,i}^A,
\label{eq:acost}\\
\rho_k^A &= \sum_i \sigma_k^A(a_{k,i}^A)\,f_1^k\,\lambda_{k,i},
\label{eq:risk}
\end{align}
where $v_{k,i}^A$ is the value of the asset targeted by action $a_{k,i}^A$,
$p_{k,i}^A$ its technical success probability, $r_{k,i}^A$ its resource cost,
and $\lambda_{k,i}$ its detection sensitivity. Crucially, the deception term
$f_1^k$ enters both $g_k^A$ (as a factor $1-f_1^k$, reducing gain when the
attacker hits a decoy) and $\rho_k^A$ (increasing the risk of being trapped).
This dual effect forces
the attacker to balance aggression against caution, producing genuinely mixed
strategies at equilibrium, as illustrated in Figure~\ref{fig:interaction}. The global attacker utility aggregates over layers
subject to the $C_1$ accessibility indicator:
\begin{equation}
U^A(\sigma^A,\sigma^D\mid\theta_A) = \sum_{k=1}^3 \beta_k^A\,
\E_{\mathbf{s}}\!\left[u_k^A\cdot\indic[\mathcal{A}_k^A(\mathbf{s})\neq\emptyset]
\right].
\label{eq:util_A_global}
\end{equation}

\begin{remark}[Structural asymmetry]
\label{rem:asymmetry}
The defender's global utility $U^D$ is an expectation over attacker
\emph{types} (via the belief $\mu_k^t$), whereas the attacker's utility $U^A$ is
an expectation over system \emph{states} (which nodes are decoys). This double
incomplete information is the precise source of the game's non-zero-sum
structure: the resilience, energy-budget, and coherence terms in $U^D$ have no
counterpart in $U^A$, so $u_k^D \neq -u_k^A$ in general. For the dual-domain
game of~\cite{asghar2025} the defender and attacker payoffs differ only by a
column-constant term, so the game is \emph{strategically equivalent} to a
zero-sum game and is solvable by a linear program. Here that property fails: the
constraint-driven penalties in $U^D$ are not column-constant, the bimatrix
$M_k^D + M_k^A$  is not rank-one, and no zero-sum reduction exists. The
restricted game must therefore be solved as a genuine non-zero-sum bimatrix
game via the LCP (Section~\ref{sec:algo}).
\end{remark}

\begin{figure}[t]
\centering
\includegraphics[width=\columnwidth]{fig_interaction.png}
\caption{Per-turn interaction between the two players. Both act simultaneously
on the deception effectiveness $f_1^k$: the defender pushes it up by deploying
more and better decoys, the attacker pushes it down through detection effort.
Because $f_1^k$ both reduces the attacker's gain and raises its risk, neither
player can play a pure strategy, yielding a mixed Nash equilibrium computed by
the LCP. The resulting state transition updates resilience $\Omega(t)$, and the
observed attacker action feeds the Bayesian belief update for the next turn.}
\label{fig:interaction}
\end{figure}

% ══════════════════════════════════════════════════════════════
\section{Resilience and the R-BNE}
\label{sec:rbne}
% ══════════════════════════════════════════════════════════════

\subsection{Resilience Function}

Following the resistance--absorption--recovery decomposition common in
cyber-physical resilience analysis~\cite{henry2012,linkov2014,zhu2015}, the local
resilience at layer $k$ is a weighted sum of three components:
\begin{equation}
\Omega_k(t) = w_1 \cdot \mathrm{res}_k(t) + w_2 \cdot \mathrm{abs}_k(t) +
w_3 \cdot \mathrm{rec}_k(t),
\end{equation}
defined as
\begin{align}
\mathrm{res}_k(t) &= \indic\!\left[\,1-\delta_k(t) \geq \rho_k\,\right], \\
\mathrm{abs}_k(t) &= \max\!\left(0,\; 1-\delta_k(t)\right), \\
\mathrm{rec}_k(t) &= \indic\!\left[\,\delta_k(t)\,|\mathcal{N}_k|\,\tau_k
\leq T_k^{\max}\,\right],
\end{align}
where $\rho_k$ is the resistance threshold, $\tau_k$ the per-node recovery time,
and $T_k^{\max}$ the maximum acceptable recovery time.
Resistance indicates whether the uncompromised fraction exceeds $\rho_k$;
absorption measures the remaining healthy fraction; and recovery
indicates whether the compromised nodes can be restored within the acceptable
time budget. The weights $w_1,w_2,w_3$ adapt to the global compromise level
$\bar\delta(t)=\frac13\sum_k\delta_k(t)$: resistance dominates when the system
is healthy, absorption when it is partially compromised, and recovery when
damage is extensive. The global resilience is the weighted aggregation
\begin{equation}
\Omega(t) = \sum_{k=1}^3 \beta_k^\Omega\, \Omega_k(t),
\qquad \sum_k\beta_k^\Omega = 1.
\end{equation}

\subsection{Solution Concept}

\begin{definition}[Admissible policy set]
\label{def:admissible}
The set $\Pi^D_\Omega$ of admissible defender policies comprises all
belief-dependent behavioral policies $\sigma^D:\Delta(\Theta_A)\to
\Delta(\mathcal{A}^D)$ that satisfy the resilience constraint
$\Omega(t)\geq\Omega_{\min}(t)$ for all $t$ and the coupling constraints
$\mathcal{C}$.
\end{definition}

\begin{definition}[R-BNE]\label{def:rbne}
A \emph{Resilience-constrained Bayesian Nash Equilibrium} is a strategy profile
$(\sigma^{D*}(\cdot),\sigma^{A*}(\cdot|\theta_A^*))$ satisfying, for all belief
states $\mu^t\in\Delta(\Theta_A)$:
\begin{align}
\text{(i)}\;\; &\sigma^{D*}(\mu^t) = \argmax_{\sigma^D\in\Pi^D_\Omega}
\E_{\theta_A\sim\mu^t}\!\left[U^D(\sigma^D(\mu^t),\sigma^{A*})\right],\\
\text{(ii)}\;\; &\sigma^{A*}(\cdot|\theta_A^*) =
\argmax_{\sigma^A\in\Delta(\mathcal{A}^A)}
U^A(\sigma^{D*},\sigma^A\mid\theta_A^*),\\
\text{(iii)}\;\; &\Omega(t)\geq\Omega_{\min}(t)\quad\forall t\geq 0.
\end{align}
\end{definition}

The R-BNE differs from a standard Bayesian Nash equilibrium~\cite{nash1951} in
three respects.
First, the defender strategy $\sigma^{D*}(\mu^t)$ is a belief-dependent
behavioral policy, mapping the current belief to a mixed action rather than
being a fixed distribution. Second, the feasible set $\Pi^D_\Omega$ imposes a
hard resilience constraint at every time step. Third, condition~(ii) uses
the attacker's true utility $U^A$, not the defender's
constraint-penalized objective, ensuring that the attacker's rationality is
modeled independently of the defender's internal penalties.

\begin{proposition}[Existence]
\label{prop:existence}
If (E1) the admissible set $\Pi^D_\Omega$ is non-empty, (E2) each $u_k^D$ is
continuous and concave in $\sigma_k^D$, and (E3) the action simplices are
compact and convex, then the MD-CDBG admits at least one R-BNE.
\end{proposition}

% ══════════════════════════════════════════════════════════════
\section{R-BNE-DO Algorithm}
\label{sec:algo}
% ══════════════════════════════════════════════════════════════

\subsection{Lagrangian Decomposition}

The coupling constraints $C_2$ and $C_3$ are dualized via non-negative Lagrange
multipliers $\nu_k$ (semantic coherence at layer $k$) and $\lambda$ (global
budget). For fixed multipliers, the Lagrangian separates into three independent
layer sub-problems:
\begin{equation}
\mathcal{L} = \sum_{k=1}^3 \mathcal{L}_k(\sigma_k^D,\sigma_k^A,\lambda,\nu_k),
\label{eq:lagrangian}
\end{equation}
where each local Lagrangian augments the belief-weighted defender utility with
the budget and coherence penalties:
\begin{align}
\mathcal{L}_k &= \beta_k\,\E_{\theta_A\sim\mu_k^t}\!\left[u_k^D\right]
- \lambda f_4^k
- \nu_k\max\!\left(0, \delta_{\min}-\Coh_k\right).
\end{align}
This separation is what enables the three sub-games to be solved in parallel.
The dual problem
$\min_{\lambda,\boldsymbol{\nu}\geq0}\max_{\sigma^D\in\Pi^D_\Omega,\sigma^A}
\mathcal{L}$ is solved by dual gradient ascent: the multipliers are updated in
the direction of their respective constraint violations,
\begin{align}
\lambda^{t+1} &= \left[\lambda^t + \eta_t\Big(\textstyle\sum_k f_4^k - B\Big)
\right]^+,\\
\nu_k^{t+1} &= \left[\nu_k^t + \eta_t\big(\delta_{\min}-\Coh_k\big)\right]^+,
\end{align}
with step $\eta_t=\eta_0/\sqrt{t}$ satisfying the Robbins--Monro conditions
$\sum_t\eta_t=\infty$ and $\sum_t\eta_t^2<\infty$. The multipliers act as
adaptive shadow prices: $\lambda$ rises when the budget is exceeded, $\nu_k$
rises when inter-layer coherence falls below $\delta_{\min}$.

\subsection{Bayesian Double Oracle}

The inner oracle (Algorithm~\ref{alg:bayes_do}) computes the local Bayesian Nash
equilibrium of sub-game $\mathcal{G}_k$ for fixed multipliers. Because the
MD-CDBG is non-zero-sum (Remark~\ref{rem:asymmetry}), the restricted-game
equilibrium cannot be computed by a single minimax linear program; instead, we
use a Linear Complementarity Program~\cite{lemke1964} with two distinct payoff
matrices. Let $e_i$ denote the pure defender action $a_{k,i}^D$ and $e_j$ the
pure attacker action $a_{k,j}^A$. The defender matrix $M_k^D$ encodes the
belief-weighted Lagrangian payoff, and the attacker matrix $M_k^A$ encodes the
attacker's true utility:
\begin{align}
M_k^D[i,j] &= \E_{\theta_A\sim\mu_k^t}\!\bigl[u_k^D(e_i,e_j\mid\theta_A)\bigr]
- \lambda f_4^k(e_i) \notag\\
&\quad - \nu_k\max\!\bigl(0,\,\delta_{\min}-\Coh_k(e_i,\sigma_{k+1}^D)\bigr),
\label{eq:MD}\\[2pt]
M_k^A[i,j] &= u_k^A(e_i,e_j\mid\theta_A^*).
\label{eq:MA}
\end{align}
Two points distinguish them. First, $M_k^D$ is an expectation over the
belief $\mu_k^t$ (the defender does not know the type), whereas $M_k^A$
uses the attacker's true type $\theta_A^*$. Second, the budget and
coherence penalties ($\lambda,\nu_k$) appear only in $M_k^D$: they are
the defender's internal constraints and do not affect the attacker's genuine
incentives. Consequently $M_k^D[i,j]+M_k^A[i,j]\neq 0$ in general, which is
precisely why a non-zero-sum LCP is required rather than a single minimax LP.
The defender's best response (Step~2) is computed in expectation over the
belief, whereas the attacker's best response (Step~3) uses its true utility,
never the Lagrangian.

\begin{algorithm}[t]
\caption{Bayes-DO: Bayesian Double Oracle for $\mathcal{G}_k$}
\label{alg:bayes_do}
\begin{algorithmic}[1]
\REQUIRE $\mathcal{G}_k$, multipliers $(\lambda,\nu_k)$, belief
$\mu_k^t(\theta_A)$
\ENSURE Local BNE $(\sigma_k^{D*},\sigma_k^{A*})$
\STATE Initialize $R_D\gets\{a_{k,1}^D\}$, $R_A\gets\{a_{k,1}^A\}$
\REPEAT
\STATE $(\tilde\sigma_k^D,\tilde\sigma_k^A)\gets\mathrm{NE}\big(M_k^D|_{R_D
\times R_A}, M_k^A|_{R_D\times R_A}\big)$ via LCP
\STATE $a_k^{D*}\gets\argmax_a \E_{\theta_A\sim\mu_k^t}
[\mathcal{L}_k(a,\tilde\sigma_k^A)]$
\STATE $a_k^{A*}\gets\argmax_a u_k^A(\tilde\sigma_k^D,a\mid\theta_A^*)$
\IF{$a_k^{D*}\in R_D$ \AND $a_k^{A*}\in R_A$}
\STATE \textbf{return} $(\tilde\sigma_k^D,\tilde\sigma_k^A)$
\COMMENT{exact equilibrium found}
\ELSE
\STATE $R_D\gets R_D\cup\{a_k^{D*}\}$;\;\; $R_A\gets R_A\cup\{a_k^{A*}\}$
\ENDIF
\UNTIL{convergence}
\end{algorithmic}
\end{algorithm}

\begin{remark}[Exact termination]
Bayes-DO is a combinatorial algorithm with exact termination: the expansion
check at Step~6 returns true precisely when no profitable deviation exists
outside the restricted sets, which is the definition of a Nash equilibrium of
the full game. No approximation tolerance is required, since equilibria are
typically supported on a small subset of actions.
\end{remark}

\subsection{Main Algorithm}

The outer algorithm R-BNE-DO (Algorithm~\ref{alg:rbne_do}) operates over a
horizon $T$ in four phases per step. Phase~A propagates the $C_1$ constraint,
checks resilience feasibility, and calls Bayes-DO per layer in parallel; if the
global resilience would fall below $\Omega_{\min}$, a fallback policy
prioritizing node isolation is activated to restore resilience. Phase~B samples
actions and transitions the state. Phase~C updates the per-layer belief.
Phase~D updates the Lagrange multipliers.

\begin{algorithm}[t]
\caption{R-BNE-DO: R-BNE via Bayesian Double Oracle}
\label{alg:rbne_do}
\begin{algorithmic}[1]
\REQUIRE MD-CDBG, $B$, $\Omega_{\min}$, $\delta_{\min}$, $\eta_0$,
$\varepsilon$, $T$
\STATE Initialize $\lambda^0\gets0$, $\nu_k^0\gets0$,
$\mu_k^0(\theta_A)\gets p(\theta_A)$
\FOR{$t=0,\ldots,T-1$}
\STATE \textbf{[A]} Propagate $C_1$; compute $\Omega(t)$
\FOR{$k=1,2,3$ \textbf{in parallel}}
\IF{$\Omega(t)<\Omega_{\min}$}
\STATE Activate fallback (isolation) policy
\ELSE
\STATE $(\sigma_k^{D,t},\sigma_k^{A,t})\gets$
Bayes-DO$(\mathcal{G}_k,\lambda^t,\nu_k^t,\mu_k^t)$
\ENDIF
\ENDFOR
\STATE \textbf{[B]} Sample $a^A(t),a^D(t)$; transition $\mathbf{s}(t{+}1)$
\STATE \textbf{[C]} Update belief $\mu_k^{t+1}$ via~\eqref{eq:bayes}
\STATE \textbf{[D]} Update $\lambda^{t+1},\nu_k^{t+1}$ by dual ascent
\IF{$\max(\|\Delta\sigma^D\|,|\Delta\lambda|,\max_k|\Delta\nu_k|)<\varepsilon$}
\STATE \textbf{return} $(\sigma^{D,t+1},\sigma^{A,t+1})$
\ENDIF
\ENDFOR
\end{algorithmic}
\end{algorithm}

\subsection{Convergence Analysis}

\begin{theorem}[Convergence]\label{thm:conv}
Under the conditions of Proposition~\ref{prop:existence}, with Robbins--Monro
steps and assuming (A1) uniqueness of the restricted Nash equilibrium at each
Bayes-DO iteration (or a fixed selection rule), and (A2) local concavity of
$\mathcal{L}_k$ in $\sigma_k^D$, the R-BNE-DO algorithm converges to an R-BNE.
\end{theorem}

\begin{proof}[Proof sketch]
Convergence combines three arguments operating at the three nested levels.
(i) Inner level: for fixed multipliers, Bayes-DO terminates in finitely
many iterations at the exact restricted Nash equilibrium, computed by the
LCP~\cite{lemke1964}; assumption (A1) prevents cycling. (ii) Outer level:
dual gradient ascent with diminishing steps converges to the saddle point of the
Lagrangian~\eqref{eq:lagrangian} under Slater's strong-duality condition, which
holds because $\Pi^D_\Omega\neq\emptyset$; at convergence the complementary
slackness conditions hold, $\lambda^*>0$ only if $\sum_k f_4^k=B$, and
$\nu_k^*>0$ only if $\Coh_k=\delta_{\min}$. (iii) Belief level: by
Lemma~\ref{lem:belief}, $\mu_k^t(\theta_A^*)\to1$ almost surely, so the
defender's best response converges to the one under perfect type knowledge,
tightening the bound on the iterations needed.
\end{proof}

\begin{proposition}[Per-call complexity]
\label{prop:complexity}
Let $r_D=|R_D|$ and $r_A=|R_A|$ be the restricted-set sizes at termination.
The LCP solved by support enumeration~\cite{porter2008} costs
$\mathcal{O}(2^{r_D}2^{r_A}(r_D+r_A)^3)$. In the MD-CDBG, $r_D\leq5$ and
$r_A\leq6$, so each call costs at most $\approx2.7\times10^6$ elementary
operations, and the number of outer iterations to reach precision $\varepsilon$
is $\mathcal{O}(1/\varepsilon^2)$.
\end{proposition}

% ══════════════════════════════════════════════════════════════
\section{Experimental Evaluation}
\label{sec:results}
% ══════════════════════════════════════════════════════════════

\subsection{Implementation}

The MD-CDBG and the R-BNE-DO algorithm were implemented from scratch in
Python~3 using \texttt{NumPy} for the numerical core and \texttt{SciPy}
for linear programming. The custom simulator faithfully mirrors
Algorithms~\ref{alg:bayes_do} and~\ref{alg:rbne_do}: it maintains a per-layer
state and belief, runs the four phases (feasibility check, Bayes-DO,
transition, belief and multiplier updates) at each of the $T$ time steps, and
records the per-step utilities, resilience, and beliefs. The non-zero-sum local
equilibria are computed by a dedicated Linear Complementarity Program
solver based on support enumeration~\cite{porter2008}, which returns the exact
mixed Nash equilibrium of each restricted bimatrix game; the zero-sum baseline
instead calls a minimax linear program through \texttt{SciPy}'s
\texttt{linprog}. Each experiment runs on a standard CPU in a few minutes; no
GPU is required. To support reproducibility, the complete source code and a
ready-to-run notebook are made available.\footnote{Source code and Colab
notebook available upon request / at the project repository.}

\subsection{Setup}

We simulate a three-layer IoT system with $|\mathcal{N}_k|=\{15,8,4\}$ real
nodes and $|\hat{\mathcal{N}}_k|=\{10,5,3\}$ decoys, over a horizon $T=50$ and
$10$ independent replications per scenario. The three scenarios correspond to
the three attacker types: O (opportunistic, $\psi=0.20$), P (persistent,
$\psi=0.50$), and E (state-sponsored, $\psi=0.65$--$0.70$). We set
$\Omega_{\min}=0.55$, $\delta_{\min}=0.60$, budget $B=1.2$, discount
$\gamma=0.95$, and step $\eta_0=0.10$. Defender utility weights are
$\alpha=(0.55,0.32,0.07,0.06)$. Parameter ranges are calibrated from the
IoT-security literature: detection efficiencies $\psi_k(\theta_A)$ follow the
attacker-sophistication levels reported for APT honeypot
detection~\cite{wan2022}, decoy ratios reflect typical honeypot-deployment
densities, and the resilience threshold $\Omega_{\min}$ is set to a
service-continuity level consistent with critical-IoT requirements. All results
are reported with 95\% confidence intervals.

\subsection{Baselines}

We compare R-BNE-DO+LCP against three baselines spanning the spectrum from
naive to state-of-the-art, chosen so that each isolates one design element:
\begin{itemize}
\item No-Deception: the defender performs passive monitoring only and
deploys no decoys; the attacker plays a rational best response. This is the
lower bound that quantifies the value of deception itself.
\item Random: decoys are deployed uniformly at random across actions,
with no strategic reasoning. Comparing against it isolates the contribution of
optimized placement over mere presence of decoys.
\item Asghar-2025~\cite{asghar2025}: the state-of-the-art GameSec~2025
double-oracle deception model, ported to our IoT setting: a greedy submodular
defender oracle with restricted games solved by a linear program under the
authors' strategic zero-sum equivalence, with no Bayesian belief and no
resilience constraint. It shares the double-oracle backbone with our method, so
comparing against it isolates the two genuine differences: the non-zero-sum
(LCP) treatment and the resilience constraint.
\end{itemize}
This design lets each comparison attribute the observed gain to a specific
component of the MD-CDBG rather than to the framework as a whole.

\subsection{Metrics}

We report three metrics. The first two are the discounted total
utilities accumulated over the horizon, where the discount factor
$\gamma\in(0,1)$ weights early time steps more heavily, reflecting the greater
value of preventing early compromise:
\begin{equation}
U^D = \sum_{t=0}^{T-1}\gamma^t\,U^D(\sigma^{D,t},\mu^t),
\qquad
U^A = \sum_{t=0}^{T-1}\gamma^t\,U^A(\sigma^{A,t}\mid\theta_A^*).
\label{eq:discounted}
\end{equation}
A higher $U^D$ is better for the defender; for $U^A$, a value closer to zero
indicates a more rational, less self-defeating attacker model (a strongly
negative $U^A$ means the attacker is forced into moves that harm itself, which
signals an unrealistic adversary model rather than a defender success). The
third metric is the violation rate
\begin{equation}
\mathrm{VC} = \frac{1}{T}\sum_{t=0}^{T-1}\indic\!\left[\Omega(t)<\Omega_{\min}
\right],
\label{eq:vc}
\end{equation}
the fraction of time steps during which the resilience constraint is breached.
A lower VC means the system maintains service continuity more consistently and
is the primary indicator of the defender's true objective. When comparing two
methods we report the relative violation reduction
$$\Delta_{\mathrm{VC}} = (\mathrm{VC}_{\text{base}}-\mathrm{VC}_{\text{ours}})
/\,\mathrm{VC}_{\text{base}}$$ i.e. the fractional decrease in violation rate of
our method with respect to a baseline; all reported metrics are averaged over
$10$ independent replications.

\subsection{Results}

\begin{figure}[t]
\centering
\includegraphics[width=\columnwidth]{fig1.png}
\caption{Defender utility $U^D$ and absolute attacker utility $|U^A|$ across
the three scenarios (10 runs, 95\% CI). The MD-CDBG induces a far more rational
attacker ($|U^A|\approx0.6$) than the zero-sum proxy of
Asghar-2025 ($|U^A|\approx2.0$).}
\label{fig:utility}
\end{figure}

\begin{figure}[t]
\centering
\includegraphics[width=\columnwidth]{fig2.png}
\caption{Resilience-constraint violation rate VC (10 runs, 95\% CI).
R-BNE-DO+LCP reduces violations by 65--72\% relative to Asghar-2025 across all
scenarios.}
\label{fig:vc}
\end{figure}

\begin{table}[t]
\centering
\caption{Four-method comparison (mean over 10 runs). Lower VC is better; less
negative $U^A$ indicates a more rational attacker.}
\label{tab:results}
\begin{tabular}{llccc}
\toprule
\textbf{Sc} & \textbf{Method} & $U^D$ & $U^A$ & VC \\
\midrule
\multirow{4}{*}{O}
& No-Deception   & 3.51 & $-0.36$ & 0.76 \\
& Random         & 5.47 & $-0.27$ & 0.30 \\
& Asghar-2025    & 5.25 & $-1.77$ & 0.60 \\
& \textbf{R-BNE-DO+LCP} & 5.62 & $-0.68$ & \textbf{0.18} \\
\midrule
\multirow{4}{*}{P}
& No-Deception   & 3.50 & $-0.35$ & 0.74 \\
& Random         & 5.24 & $-0.30$ & 0.28 \\
& Asghar-2025    & 5.25 & $-1.77$ & 0.60 \\
& \textbf{R-BNE-DO+LCP} & \textbf{5.56} & $-0.63$ & \textbf{0.21} \\
\midrule
\multirow{4}{*}{E}
& No-Deception   & 3.71 & $-0.32$ & 0.71 \\
& Random         & 5.26 & $-0.27$ & 0.31 \\
& Asghar-2025    & 5.19 & $-1.78$ & 0.64 \\
& \textbf{R-BNE-DO+LCP} & \textbf{5.79} & $-0.56$ & \textbf{0.18} \\
\bottomrule
\end{tabular}
\end{table}

\begin{figure*}[t]
\centering
\includegraphics[width=\textwidth]{fig3.png}
\caption{Global resilience $\Omega(t)$ over time (10 runs, shaded $\pm$std).
R-BNE-DO+LCP maintains $\Omega(t)$ above $\Omega_{\min}$ far more consistently
than the baselines.}
\label{fig:resilience}
\end{figure*}

Table~\ref{tab:results} summarizes the four-method comparison, and
Figures~\ref{fig:utility}--\ref{fig:resilience} present the same data
graphically.

Resilience (Fig.~\ref{fig:vc}). This is the headline result.
R-BNE-DO+LCP achieves the lowest violation rate in every scenario
(VC~$=0.18$--$0.21$). Applying the relative violation reduction
$\Delta_{\mathrm{VC}}$ defined above, this is a $69\%$, $65\%$, and $72\%$
decrease relative to Asghar-2025 in scenarios O, P, and E respectively (e.g.
$(0.60-0.18)/0.60\approx0.69$ for O), and up to $76\%$ relative to
No-Deception. The full ordering is consistent across scenarios:
R-BNE-DO~$<$~Random~$<$~Asghar-2025~$<$~No-Deception. That
Random outperforms the more sophisticated Asghar-2025 on
resilience is itself revealing: methods that optimize a non-resilience objective
can actively steer the system into low-resilience states, whereas random
placement at least spreads risk.

Defender utility (Fig.~\ref{fig:utility}a). R-BNE-DO+LCP dominates in
scenarios P ($5.56$ vs $5.25$) and E ($5.79$ vs $5.19$), and is statistically
indistinguishable from the best baseline in scenario O (overlapping 95\%
confidence intervals). The interpretation is important: against the weakest
attacker (O), aggressive greedy placement and equilibrium placement perform
comparably on raw utility, but our method still secures the system far better
(VC $0.18$ vs $0.60$). In other words, the defender obtains equal or higher
utility and substantially better resilience, it does not trade one for
the other.

Attacker rationality (Fig.~\ref{fig:utility}b). The zero-sum proxy of
Asghar-2025 drives the attacker utility to $\approx-2.0$, i.e. an adversary
systematically forced into self-harming moves. Our non-zero-sum treatment yields
$U^A\approx-0.6$, a rational attacker that plays to its genuine incentives. This
gap is not a defender win, it signals that the zero-sum model misrepresents
the threat, which can lead to defenses tuned against an unrealistic adversary.

Temporal behavior (Fig.~\ref{fig:resilience}). The resilience
trajectories show that R-BNE-DO+LCP stays above $\Omega_{\min}$ most of the
time and rebounds quickly after compromise spikes, a direct effect of the
fallback isolation policy triggered by the feasibility check. The baselines, by
contrast, drift below the threshold and remain there, never recovering.

\subsection{Discussion}

Beyond the per-figure reading above, two higher-level lessons emerge. First, the
comparison design (Section~\ref{sec:results}) lets us attribute the resilience
gain to specific causes: the gap to Asghar-2025 isolates the value of the
non-zero-sum (LCP) treatment and the resilience constraint, while the gap to
Random isolates the value of strategic, equilibrium-based placement over the
mere presence of decoys. Both gaps are large, confirming that neither component
alone suffices. Second, the consistency of the ordering across three very
different attacker profiles suggests the advantage is structural rather than
tuned to a particular adversary: embedding resilience as a hard constraint
protects service continuity regardless of who is attacking.

══════════════════════════
\section{Conclusion}
\label{sec:conclusion}
% ══════════════════════════════════════════════════════════════

We presented the MD-CDBG, a multi-domain cyber-deception Bayesian game that
couples the three IoT layers under formal resilience constraints, together with
the R-BNE solution concept and the R-BNE-DO algorithm. By correctly modeling the
non-zero-sum structure via an LCP oracle and embedding resilience as a hard
constraint, the framework reduces resilience violations by 65--72\% relative to
the state-of-the-art GameSec~2025 approach while inducing rational attacker
behavior.

Several directions extend this work; while the present results are established
in simulation, the framework is designed to carry over to deployed settings. (i) Learning-based oracle: the
greedy and equilibrium oracles could be complemented by a trained deep detector
(e.g., a graph attention network or transfer autoencoder) acting as an
additional best-response oracle within the double-oracle loop, combining the
forward-looking guarantees of game theory with the pattern-recognition strength
of machine learning. (ii) Symmetric information: a symmetric extension
would model the attacker's own belief over the defender
type, turning the game into a two-sided incomplete-information interaction.
(iii) Richer deception techniques: the defender's action repertoire can be
extended with additional network-layer mechanisms such as software-defined
network reconfiguration and deceptive routing or address-shuffling, each adding
new defender actions while preserving the multi-domain coupling. (iv) Scalability: approximate equilibrium solvers
would let the model scale to IoT deployments with thousands of heterogeneous
nodes. Together, these directions point toward resilient-by-design
cyber-deception, where security and service continuity are optimized jointly
rather than traded off against each other.

\section*{Acknowledgment}
This work was supported by the French Government Scholarship (BGF) and conducted
within the cotutelle program between Universit\'e Paris-Saclay (Laboratoire
DAVID) and Universit\'e de Ngaound\'er\'e.

\bibliographystyle{IEEEtran}
\bibliography{references}

\end{document}
